\section{Taller resuelto: Una cizalladura de viento tiene aceleración}

\noindent\textbf{Nivel:} Intermedio.

\apartado{1. Enunciado}{Un campo estacionario tiene viento norte U(h)=0.003h m/s. Un avión asciende a 4 m/s y mantiene velocidad aérea norte constante. Determine la aceleración norte de su velocidad terrestre y la fuerza adicional que supone para 2600 kg.}

\apartado{2. Datos}{$\partial U/\partial h=0.003\SIu{s^{-1}}$, $\dot h=4\SIu{m/s}$.}

\apartado{3. Incógnitas}{Las magnitudes y comprobaciones solicitadas en el caso.}

\apartado{4. Sistema de referencia}{Componente norte local y altitud positiva hacia arriba; referencia plana.}

\apartado{5. Hipótesis físicas}{Se considera el gradiente indicado y se omiten otros cambios de viento y actitud.}

\apartado{6. Modelo matemático}{Una cizalladura de viento tiene aceleración}

\apartado{7. Ecuaciones simbólicas}{Se parte de los balances y definiciones siguientes.}

\begin{equation}v_{g,N}=v_{a,N}+U(h)\end{equation}

\begin{equation}\dot v_{g,N}=\dot v_{a,N}+(\dd U/\dd h)\dot h\end{equation}

\begin{equation}\Delta F_N=m(\dd U/\dd h)\dot h\end{equation}

\apartado{8. Desarrollo algebraico}{Las operaciones se realizan en este orden, conservando las condiciones del modelo.}

\begin{enumerate}

\item El viento no cambia explícitamente con el reloj, pero el avión recorre su gradiente.

\item La derivada de viento a lo largo de la trayectoria usa hdot.

\item El producto por masa da el término inercial adicional; no se lo añade como fuerza aerodinámica ficticia.

\end{enumerate}

\apartado{9. Sustitución numérica}{Se conservan los resultados intermedios sin redondear y se muestran todas las operaciones escalares disponibles.}

\begin{enumerate}

\item Aceleración de viento.

\[\fitmath{y_{1}=0.003\cdot 4\approx 0.012\,\mathrm{m/s^{2}}}\]

\item Término de fuerza.

\[\fitmath{y_{2}=2600\cdot 0.003\cdot 4\approx 31.2\,\mathrm{N}}\]

\end{enumerate}


\apartado{10. Resultados con unidades}{Los resultados pertenecen al modelo y al intervalo de datos declarados.}

\begin{center}\begin{tabular}{lr}\toprule Magnitud&Resultado\\\midrule

Aceleración de viento & $0.012\,\mathrm{m\,s^{-2}}$\\

Término de fuerza & $31.2\,\mathrm N$\\

\bottomrule\end{tabular}\end{center}

\apartado{11. Interpretación}{Mantener velocidad relativa al aire constante en cizalladura no implica velocidad terrestre constante.}

\apartado{12. Comprobación independiente y alcance}{Tomar dos instantes separados 10 s: h aumenta 40 m y U aumenta 0.12 m/s; dividir por tiempo recupera aceleración.}

\apartado{13. Limitaciones}{La pregunta es local y escalar; una cizalladura real puede alterar incidencia, cargas y control de modo adicional.}

\par\smallskip\noindent{\footnotesize\textit{Datos y solución originales. Las sustituciones registradas se reproducen con SymPy y Studium; eso no establece validación experimental ni revisión experta.}}
